%%%% ijcai26.tex

\typeout{IJCAI-ECAI 26 Instructions for Authors}

% These are the instructions for authors for IJCAI-ECAI 26.

\documentclass{article}
\pdfpagewidth=8.5in
\pdfpageheight=11in

% The file ijcai26.sty is a copy from ijcai22.sty
% The file ijcai22.sty is NOT the same as previous years'
\usepackage{ijcai26}

% Use the postscript times font!
\usepackage{times}
\usepackage{soul}
\usepackage{url}
\usepackage[hidelinks]{hyperref}
\usepackage[utf8]{inputenc}
\usepackage[small]{caption}
\usepackage{graphicx}
\usepackage{amsmath}
\usepackage{amsthm}
\usepackage{booktabs}
\usepackage{algorithm}
\usepackage{algorithmic}
\usepackage[switch*]{lineno}
% lineno's first two-column page is recorded as page 0 before its first
% column marker is written. Seed that marker for the left column so both
% columns place their numbers in the outer margins.
\makeatletter
\def\@LN@column{2}
\makeatother
\usepackage{xcolor}
\usepackage{multirow}
\usepackage{float}
\usepackage{subcaption}
\usepackage{amsfonts}
\usepackage{adjustbox}
\usepackage{wrapfig}

% Comment out this line in the camera-ready submission

\urlstyle{same}

% the following package is optional:
%\usepackage{latexsym}

% See https://www.overleaf.com/learn/latex/theorems_and_proofs
% for a nice explanation of how to define new theorems, but keep
% in mind that the amsthm package is already included in this
% template and that you must *not* alter the styling.
\newtheorem{example}{Example}
\newtheorem{theorem}{Theorem}

% Following comment is from ijcai97-submit.tex:
% The preparation of these files was supported by Schlumberger Palo Alto
% Research, AT\&T Bell Laboratories, and Morgan Kaufmann Publishers.
% Shirley Jowell, of Morgan Kaufmann Publishers, and Peter F.
% Patel-Schneider, of AT\&T Bell Laboratories collaborated on their
% preparation.

% These instructions can be modified and used in other conferences as long
% as credit to the authors and supporting agencies is retained, this notice
% is not changed, and further modification or reuse is not restricted.
% Neither Shirley Jowell nor Peter F. Patel-Schneider can be listed as
% contacts for providing assistance without their prior permission.

% To use for other conferences, change references to files and the
% conference appropriate and use other authors, contacts, publishers, and
% organizations.
% Also change the deadline and address for returning papers and the length and
% page charge instructions.
% Put where the files are available in the appropriate places.

% PDF Info Is REQUIRED.

% Please leave this \pdfinfo block untouched both for the submission and
% Camera Ready Copy. Do not include Title and Author information in the pdfinfo section
\pdfinfo{ /TemplateVersion (IJCAI.2026.0) }

\title{Bridging Evolution and Learning: Attention-Guided Co-Design of Modular Soft
Robots}

% Single author syntax
\author{ Author Name \affiliations Affiliation \emails email@example.com }

% Multiple author syntax (remove the single-author syntax above and the \iffalse ... \fi here)
\iffalse
\author{ First Author$^{1}$ \and Second Author$^{2}$\and Third Author$^{2,3}$\And
Fourth Author$^{4}$\\
\affiliations $^{1}$First Affiliation\\
$^{2}$Second Affiliation\\
$^{3}$Third Affiliation\\
$^{4}$Fourth Affiliation\\
\emails \{first, second\}@example.com, third@other.example.com, fourth@example.com
}
\fi
\newcommand{\red}[1]{\textcolor{red}{#1}}
\begin{document}
    \maketitle
    \linenumbers
    \begin{abstract}
        Jointly optimizing the morphology and controller of a voxel-based soft
        robot is expensive because evaluating each candidate body requires
        controller training. Under a fixed training budget, the search must
        explore diverse morphologies while identifying candidates that warrant
        further optimization. We propose Attention-Guided Co-Design (AGCD), which
        combines staged controller training with attention-guided structural
        variation. A self-attention controller produces voxel-level scores from
        rollouts; crossover and mutation use these scores to favor the retention
        of highly attended voxels while varying other regions. We use attention
        as a search heuristic rather than a measure of causal voxel importance.
        AGCD also ranks candidates by observed fitness and training maturity,
        a discrete measure of cumulative controller-training effort, through
        non-dominated sorting. This allows promising partially trained robots
        to remain in the search. An elite-training phase subsequently completes
        the preset per-candidate training budget for selected robots.
        Across 27 Evolution Gym
        tasks under matched training budgets, AGCD achieves the highest final
        mean reward on 23 tasks compared with a genetic algorithm, Bayesian
        optimization, and an action-inheritance evolutionary method.
    \end{abstract}

    \section{Introduction}
    Robot morphology and control are tightly coupled: a change in body structure
    alters the dynamics that a controller must learn, and the apparent quality of
    a body depends on how thoroughly its controller has been trained. This
    interaction is fundamental to embodied intelligence~\cite{liu2025embodied,pfeifer2006morphological}
    and has motivated automated searches over bodies and behaviors since early
    evolutionary robotics studies~\cite{sims1994evolving,lipson2000automatic,bongard2003evolving}.
    The challenge is especially acute for voxel-based soft robots. Even a small
    grid admits a large number of material layouts, while evaluating each layout
    requires costly controller training~\cite{bhatia2021evolution,liu2023rapidly}.
    Allocating the full preset controller-training budget to every candidate
    can therefore exhaust the search budget before enough distinct morphologies
    are explored.

    Efficient co-design under a fixed budget requires decisions about both
    evaluation and variation: which candidates should receive more controller
    training, and how should existing candidates inform the next generation?
    Population-based methods explore multiple morphologies in parallel~\cite{stanley2009hypercube,tanaka2022co,harada2024lamarckian},
    but selecting among them remains difficult when fitness estimates depend on
    training effort. Moreover, a structural change can diminish the usefulness
    of a previously trained controller~\cite{cheney2018scalable}.
    Reinforcement-learning-based approaches provide another route to joint
    design~\cite{wang2023curriculum,wang2019neural,lu2025bodygen,xu2021end}, but
    they still require training resources to be allocated across candidate
    structures. A useful co-design method should therefore manage incomplete
    evaluations and exploit information gained during controller training to
    guide the morphology search.

    We propose Attention-Guided Co-Design (AGCD) to address these two decisions.
    To guide structural variation, AGCD aggregates attention weights from
    controller rollouts into voxel-level scores. Crossover and mutation use
    these scores to favor highly attended regions while varying other parts of
    the robot. The scores serve as search heuristics, without assuming that
    attention establishes a voxel's causal importance. To allocate training,
    AGCD evaluates candidates in stages and uses Maturity-Fitness Non-Dominated
    Sorting to rank them by observed fitness and training maturity, which
    measures cumulative controller-training effort. This ranking preserves
    promising candidates that have received relatively little training, while
    a subsequent elite phase completes the preset per-candidate training budget
    for selected archived candidates. Our main contributions are:
    \begin{itemize}
        \item We propose AGCD, a co-design framework that uses information from
            controller learning to guide both morphology evolution and
            training-budget allocation for voxel-based soft robots.
        \item We develop attention-guided crossover and mutation operators that
            preferentially retain highly attended voxels while varying other
            regions, incorporating controller-derived information into offspring
            generation.
        \item We introduce a training-budget allocation strategy that combines
            staged evaluation, Maturity-Fitness Non-Dominated Sorting, and elite
            training to retain promising undertrained candidates and selectively
            complete their preset training budgets.
    \end{itemize}

    \section{Related Work}
    \subsection{Robot Co-design}
    Robot co-design aims to jointly optimize morphology and control strategies. Due
    to the close interdependence between morphology and the controllers~\cite{rosendo2017trade,ma2023comprehensive},
    morphological design often directly influences the expressive capacity of
    control strategies. Previous studies have examined how different encoding approaches
    affect the co-optimization performance~\cite{pigozzi2023morphology,bhatia2021evolution,stanley2007compositional},
    and these methods are commonly categorized as direct encoding or indirect
    encoding. Among indirect methods, the compositional pattern-producing
    network (CPPN)~\cite{stanley2007compositional} is an efficient approach able
    to generate morphologies with symmetry and repetitive structures. In contrast,
    direct encoding often uses graph representations for robot morphologies~\cite{wang2019neural},
    allowing effective modeling of relationships between modules. Beyond single-task
    co-design, this paradigm has been extended to multi-task settings~\cite{mertan2024towards},
    facilitating knowledge transfer from simple tasks to more complex ones~\cite{zhao2025cross,wang2023preco}.
    In addition, generative models were incorporated into co-design, including classifier-guided
    diffusion models to improve the efficiency of morphological evolution~\cite{wang2023diffusebot,liu2025morphology}.
    Large language models have also been explored as evolutionary search
    operators to balance optimization efficiency with solution diversity~\cite{song2025laser,fang2025robomore}.

    Age--fitness selection also addresses the challenge of evaluating newly
    modified morphologies. Cheney et al.~\cite{cheney2018scalable} define
    morphological age as the number of generations since a morphological
    mutation and combine it with fitness to allow controller readaptation.
    AGCD shares the motivation of retaining candidates whose controllers have
    had limited optimization, but records the controller-training effort
    allocated to each candidate as training maturity. This quantity supports
    staged PPO evaluation and subsequent completion of selected candidates'
    training budgets within a fixed total budget.

    % \subsection{Non-dominated Sorting Genetic Algorithm}
    % The Non-dominated Sorting Genetic Algorithm (NSGA) is a multi-objective evolutionary
    % search algorithm. Early non-dominated sorting was applied to genetic
    % algorithms~\cite{saito2022comparative}, preserving solution diversity through
    % hierarchical dominance levels combined with fitness sharing. To address the
    % shortcomings in computational efficiency, elitism, and diversity
    % preservation in the original non-dominated sorting method, NSGA-II was proposed~\cite{srinivas1994muiltiobjective}.
    % NSGA-II employs crowding distance instead of fitness sharing to maintain a
    % well-distributed set of solutions. Due to a favorable balance between convergence
    % and diversity, evolutionary methods based on non-dominated sorting have been
    % widely applied to complex optimization problems~\cite{ma2023comprehensive}. Furthermore,
    % in multi-agent reinforcement learning, such multi-objective evolutionary strategies
    % have demonstrated strong collaborative capabilities~\cite{deb2002fast,yifei2020multi,s25041001}.
    % In this work, multi-objective ranking criteria are employed to guide the search
    % process, to improve search efficiency.
    \begin{figure}[t]
        \centering
        \begin{subfigure}
            [t]{0.60\linewidth}
            \centering
            \includegraphics[width=\linewidth, keepaspectratio]{
                pdf/structure.pdf
            }
            \caption{}
            \label{fig:left}
        \end{subfigure}
        \hfill
        \begin{subfigure}
            [t]{0.32\linewidth}
            \centering
            \includegraphics[width=\linewidth, keepaspectratio]{pdf/matrix.pdf}
            \caption{}
            \label{fig:right}
        \end{subfigure}
        \caption{Example of a VSR. (a) Robot structure. (b) Corresponding
        digital encoding.}
        \label{fig:example_robot}
    \end{figure}

    \subsection{Self-Attention}
    A controller for a voxel robot must combine information from spatially
    separated modules. Self-attention offers a way to exchange information among
    voxel representations without fixing a local communication neighborhood~\cite{vaswani2017attention,zhao2025multi}.
    For an input sequence $X$, learned projections produce the query, key, and
    value vectors $Q=XW_{Q},K=XW_{K},V=XW_{V}$. Their attention weights are
    $\mathrm{softmax}\left( \frac{QK^{\top}}{\sqrt{d_{k}}}\right )$.
    In AGCD, the controller uses these weights for action generation, while the
    morphology operator aggregates them into a voxel-level variation heuristic.
    A large weight indicates where this trained policy places attention in the
    sampled rollout; it does not by itself demonstrate that a voxel is necessary
    for the resulting behavior.

    \begin{figure*}[t]
        \centering
        \includegraphics[width=\textwidth]{pdf/method_overview.pdf}
        \caption{Overview of AGCD. It consists of two phases: the Maturity
        Evolution Phase, which evolves high-quality individuals using Maturity–Fitness
        Non-Dominated Sorting and an Attention-Based Variation Operator; and the
        Elite Training Phase, which completes the preset per-candidate
        training budget for selected archived individuals. Training maturity
        records allocated controller-training effort rather than convergence.}
        \label{fig:method_overview}
    \end{figure*}

    \section{Problem Formulation}

    The search studied here assigns a material to each cell of a two-dimensional
    design grid, then evaluates the resulting robot after controller training.
    This grid representation makes structural changes discrete, but each change
    can also alter the behavior that a controller must learn. Voxelized designs
    have been used to study automated soft-robot design and morphological
    diversity~\cite{hiller2011automatic,cheney2014unshackling}. We use the
    Evolution Gym representation~\cite{bhatia2021evolution}: the five cell
    values are empty (0), rigid (1), soft (2), horizontal actuator (3), and
    vertical actuator (4), as illustrated in Figure~\ref{fig:example_robot}.
    The two actuator types change length along their specified axes; soft
    material deforms passively, and rigid material retains its shape. Neighboring
    occupied voxels transmit forces, allowing the assembled body to move.

    Let the morphology be selected from a discrete set and its controller from
    a parameterized policy family. The idealized task objective has an outer
    choice of morphology and an inner choice of controller for that morphology:
    \begin{equation}
        \begin{aligned}
            (s^{*}, \phi_{s}^{*})      & = \operatorname*{arg\,max}_{s \in \mathcal{S}}\; F\!\left(s, \phi_{s}^{*}\right), \\
            \text{s.t.}\; \phi_{s}^{*} & = \operatorname*{arg\,max}_{\phi \in \Theta}\; f\!\left(s, \phi\right).
        \end{aligned}
        \label{eq:bilevel}
    \end{equation}
    Here, $s\in\mathcal{S}$ specifies the grid of voxel types, and $\phi\in\Theta$
    parameterizes its neural controller. Without connectivity and actuator
    constraints, an $n\times n$ grid has $5^{n^2}$ possible assignments; only
    a subset forms valid robots. The inner objective $f(s,\phi)$ describes
    controller performance for a fixed body, while the outer objective
    $F(s,\phi_s^*)$ scores a body with its optimized controller. In practice,
    reinforcement learning (RL) only approximates this inner optimum with a
    finite number of simulation steps. Fitness can thus change as controller
    training proceeds. AGCD allocates this limited training budget across
    candidate morphologies instead of requiring completion of the preset
    per-candidate training budget before each structural selection decision.

    \section{Proposed Method}
    AGCD alternates between allocating controller-training effort and generating
    new robot structures. Figure~\ref{fig:method_overview} shows its Maturity
    Evolution and Elite Training phases. During evolution, partially trained
    robots receive fitness estimates and contribute attention signals to the
    archive. Maturity-Fitness Non-Dominated Sorting selects survivors, and the
    attention-based operators construct the next generation. Elite Training then
    completes the preset per-candidate training budget for selected archived
    robots. Algorithm~\ref{alg:AGCD_pseudocode}
    gives the complete procedure. Sections~\ref{subsec:multi_maturity} and
    \ref{subsec:maturity_fitness} define evaluation and ranking;
    Section~\ref{subsec:attention_variation} specifies the structural operators.
    Section~\ref{subsec:control_strategy} describes the adopted controller and
    its training strategy.

    \subsection{Training Maturity and Staged Evaluation}
    \label{subsec:multi_maturity} A robot's observed fitness depends on the
    training already invested in its controller. We define training maturity
    as the cumulative controller-training effort allocated to a candidate,
    measured in fixed training intervals. Hereafter, maturity refers to this
    measure of training effort, rather than the controller's degree of convergence
    or the number of generations the candidate has survived.

    During the Maturity Evolution phase, each evaluation allocates one fixed
    interval of controller training and increments the candidate's training
    maturity by one. During Elite Training, selected archived candidates
    receive the remaining training needed to reach the preset per-candidate
    budget. The maximum level of training maturity is 25, as examined in
    Appendix~\ref{subsec:max_maturity}. Reaching this level indicates that the
    preset budget has been exhausted; it does not imply that the controller
    has converged. Each candidate therefore has an observed fitness and a
    level of training maturity. The following ranking rule uses both quantities
    to compare performance relative to the training effort allocated.

    \subsection{Maturity-Fitness Non-Dominated Sorting}
    \label{subsec:maturity_fitness} Non-dominated sorting groups candidates by
    Pareto dominance~\cite{srinivas1994muiltiobjective,deb2002fast}. For AGCD,
    higher fitness is preferred, while lower training maturity represents
    less allocated controller-training effort. Candidate $a$ dominates
    candidate $b$ under the following rule:
    \begin{equation}
        \label{eq:mf_dominance}a\text{ dominates }b = \left\{
        \begin{array}{l}
            a_m \le b_m,                                       \\
            a_f \ge b_f,                                       \\
            \left(a_m < b_m\right) \lor \left(a_f > b_f\right)
        \end{array}
        \right.
    \end{equation}
    The subscripts $m$ and $f$ indicate training maturity and observed fitness,
    respectively. Thus, $a$ must have received no more controller training than
    $b$ and have at least as high a fitness estimate, with one comparison
    strict. This rule allows a candidate with a strong observed score and
    limited training to remain competitive with candidates that have received
    more training. Training maturity provides a measure of allocated effort;
    it does not estimate convergence or predict performance after further
    training.

    The population is partitioned into successive non-dominated fronts. We use
    the fast sorting procedure of NSGA-II~\cite{deb2002fast}, as specified in
    Algorithm~\ref{alg:mfnds}, and order candidates within a front by fitness.
    This ranking retains some inexpensive, promising evaluations for subsequent
    search and training.

    \subsection{Attention-Based Variation Operator}
    \label{subsec:attention_variation} The trained controller allocates different
    attention weights across voxel positions. AGCD uses this allocation as a
    structural search signal, not as a causal explanation of robot behavior.
    After evaluating an individual, we run its controller for a simulation
    rollout and collect attention matrices at each time step. Averaging over
    time and normalizing yields a score for each voxel. Crossover and mutation
    use these scores to decide which regions to preserve or change.
    \begin{figure}[h]
        \centering
        \includegraphics[width=\linewidth]{pdf/crossover.pdf}
        \caption{Example of the Attention-Based Crossover Operation.}
        \label{fig:crossover}
    \end{figure}
    \subsubsection{Attention-Based Crossover Operation}
    The crossover operator uses parental attention scores to prefer some voxel
    assignments over others when assembling an offspring. Let $v_{i}^{c}$ and
    $a_{i}^{c}$ denote the encoded voxel type and its attention score for parent
    $P_{i}$ ($i \in \{1, 2\}$) at position $c$. The offspring voxel is assigned
    by the following rule:
    \begin{equation}
        v_{off}^{c}=
        \begin{cases}
            v_{\operatorname*{arg\,max}\limits_{i} a_{i}^{c}}^{c} & \text{if }v_{1}^{c}, v_{2}^{c}\neq0                        \\
            v_{i}^{c}                                             & \text{if only }v_{i}^{c}\neq0, a_{i}^{c}\geq \tilde{a}_{i} \\
            mutate(c)                                             & \text{if only }v_{i}^{c}\neq0, a_{i}^{c}< \tilde{a}_{i}    \\
            0                                                     & \text{if }v_{1}^{c}= v_{2}^{c}=0
        \end{cases}
    \end{equation}
    where $\tilde{a}_{i}$ is the median attention weight of all non-empty voxels
    in parent $P_{i}$, and $mutate(c)$ is a random mutation that changes the
    voxel at position $c$ to another voxel type. After offspring generation,
    offspring undergo structural feasibility checks. If the requirements are not
    satisfied, the process is repeated until the predefined upper limit is
    reached. Figure~\ref{fig:crossover} illustrates an example of this operation,
    where pink borders represent voxels inherited by the offspring from Parent1,
    purple borders represent voxels inherited from Parent2, and green borders
    denote newly mutated voxels.
    \begin{figure}[h]
        \centering
        \includegraphics[width=0.78\linewidth]{pdf/mutation.pdf}
        \caption{Example of the Attention-Based Mutation Operation.}
        \label{fig:mutation}
    \end{figure}
    % This operation not only enhances the stability of
    % structural inheritance by preventing high-value structures from being
    % disrupted during crossover but also preserves necessary randomness in low-contribution
    % regions, ensuring diversity within the topological structure.
    \subsubsection{Attention-Based Mutation Operation}
    The attention-based mutation operation is similar to the above crossover operation,
    taking a single parent structure and its associated attention weights as
    input. Let $v^{c}$ and $a^{c}$ denote the numerically encoded voxel type and
    its corresponding attention weight at position $c$ in parent $P$. The voxel
    of the offspring at each position is updated according to the following rule:
    \begin{equation}
        v_{off}^{c}=
        \begin{cases}
            v^{c}     & \text{if }v^{c}\neq0 \text{ and }a^{c}\geq \tilde{a} \\
            mutate(c) & \text{if }v^{c}\neq0 \text{ and }a^{c}< \tilde{a}    \\
            mutate(c) & \text{if }v^{c}=0
        \end{cases}
    \end{equation}

    where $\tilde{a}$ is the median attention weight of all non-empty voxels in
    parent $P$, and $mutate(c)$ is a random mutation that changes the voxel at position
    $c$ to another voxel type. After offspring generation, structural feasibility
    checks are performed. If requirements are not met, the process retries until
    reaching the predefined upper limit. Figure~\ref{fig:mutation} illustrates
    an example of this operation, depicting that regions with high attention weights
    are preserved while low-scoring areas are mutated into different voxels.

    Both operators preferentially retain voxels with higher attention scores
    and allow changes elsewhere. Their value as search operators is assessed
    through the experimental comparisons and ablations.
    % This operation simultaneously avoids
    % damaging critical structures while inducing controlled perturbations in low-contribution
    % regions, thereby enhancing both the diversity of structural exploration and
    % overall evolutionary efficiency.
    \subsection{Control and Training Strategies}
    \label{subsec:control_strategy}

    To facilitate rapid validation of the attention-guided structural variation
    operators, we adopt the self-attention-based control strategy of
    Wang et al.~\cite{wang2023curriculum}. The controller represents the local
    observations and material information of individual voxels as a sequence.
    Self-attention captures dependencies among voxels, allowing the policy to
    coordinate their behavior and generate actions for the actuator voxels.
    In AGCD, attention weights collected during controller rollouts are
    aggregated into the voxel-level scores used by the crossover and mutation
    operators described above.

    To train the control policy, we employ Proximal
    Policy Optimization (PPO)~\cite{schulman2017proximal} as the core reinforcement
    learning framework. PPO enhances the stability of the optimization process
    by limiting the magnitude of updates between old and new policies.

    \section{Experiments}
    We evaluate AGCD through comparative experiments and ablation studies,
    addressing two questions: (1) Can attention effectively guide morphology
    search in co-design? (2) How do attention-guided variation and allocation
    based on training maturity contribute to co-design performance?
    \subsection{Experimental Settings}
    All experiments were implemented on the Evolution Gym~\cite{bhatia2021evolution}.
    The maximum iteration count per individual for each task is referenced in Table~\ref{tab:env_training_iters}.
    We set the design space of the robot to 5x5, defined the maximum training iterations
    $t$ per individual for each task, and set the total iterations to $100t$ for
    easy tasks, $200t$ for medium tasks, and $300t$ for hard tasks. To comprehensively
    evaluate the performance of AGCD, we validated our method on 27 task
    environments, comprising 9 easy tasks, 9 medium tasks, and 9 hard tasks.
    These tasks include Walking, Object Manipulation, Climbing, and Locomotion.
    To better compare differences with other methods, we designed 3 baseline categories
    as follows:
    \begin{itemize}
        \item \textbf{GA}: Genetic Algorithm is a classic evolutionary computation
            method that simulates natural selection, inheritance, and mutation mechanisms
            to select and mutate candidate solutions, iteratively generating
            superior individuals.

        \item \textbf{BO}: Bayesian Optimization is a surrogate-assisted algorithm
            used to solve expensive black-box optimization problems~\cite{wang2023recent}.
            It typically employs Gaussian processes to approximate the true
            objective function, thereby reducing computational costs associated with
            expensive fitness evaluations.

        \item \textbf{AIEA}: Act\_Inh-based EA is a fast fitness evaluation algorithm~\cite{liu2023rapidly}.
            Candidate robots inherit actions generated by well-trained
            controllers, combine them with error estimation, and obtain
            approximate task performance.
    \end{itemize}
    AGCD adopts the $(\mu+\lambda)$ evolutionary strategy~\cite{beyer2002evolution},
    where $\mu=10$ and $\lambda=10$. To maintain consistency across experiments,
    the population size for all other methods was set to 20, and PPO parameter
    settings remained identical. For details, please refer to Appendix~\ref{subsec:hyperparameters}.

    \subsection{Results and Analysis}
    \begin{figure}[t]
        \centering
        \includegraphics[width=\linewidth]{pdf/compare_part.pdf}
        \caption{Mean and standard deviation of reward curves for selected tasks
        with AGCD, AIEA, GA, and BO. Each curve is averaged over five
        independent runs.}
        \label{fig:compare_part}
    \end{figure}
    \begin{table}[h]
        \centering
        \begin{adjustbox}
            {max width=\linewidth}
            \begin{tabular}{ccccc}
                \toprule \textbf{Task} & \textbf{AGCD(Ours)}  & \textbf{AIEA} & \textbf{GA} & \textbf{BO} \\
                \midrule Walker-v0     & *\textbf{10.52±0.25} & *10.35±0.43   & 9.40±0.87   & 7.58±0.39   \\
                Flipper-v0             & \textbf{5.09±0.30}   & 4.67±0.23     & 4.09±0.68   & 2.76±0.53   \\
                Pusher-v1              & \textbf{0.60±0.02}   & 0.19±0.01     & 0.16±0.04   & 0.10±0.01   \\
                Thrower-v0             & \textbf{2.21±0.16}   & 0.94±0.06     & 0.83±0.06   & 0.84±0.02   \\
                Hurdler-v0             & \textbf{3.05±0.21}   & 1.49±0.10     & 1.28±0.08   & 1.04±0.25   \\
                GapJumper-v0           & \textbf{8.26±0.44}   & 3.53±0.19     & 2.21±0.31   & 3.08±0.37   \\
                \bottomrule
            \end{tabular}
        \end{adjustbox}
        \caption{Mean and standard error of all results across selected tasks. Best
        results are highlighted in bold.* indicates that the result reached the upper
        bound.}
        \label{tab:part_performance}
    \end{table}
    \begin{figure}[h]
        \centering
        \includegraphics[width=0.5\textwidth]{pdf/visual_GapJumper-v0.pdf}
        \caption{Visualization results of AGCD, AIEA, GA, and BO on the GapJumper-v0
        task. The figure shows changes in robot position and motion over time
        for each method. Snapshots are arranged from left to right in chronological
        order. The evolved AGCD robot demonstrated crab-like lateral running
        with the ability to jump across gaps.}
        \label{fig:visual}
    \end{figure}

    In Figure~\ref{fig:compare_part} and Table~\ref{tab:part_performance}, we
    present the reward curves and final mean values of AGCD, AIEA, GA, and BO on
    selected representative tasks. The results indicate that on easy tasks, our method
    performed slightly better than the other methods, with small performance
    differences. However, for medium and hard tasks, our method clearly achieved
    a superior performance compared to the other methods. (For space consideration,
    the reward curves and final fitness values for all tasks are provided in
    Appendix~\ref{subsubsec:full_results}). This phenomenon mainly results from
    the low challenge of easy tasks, which makes it harder to observe differences
    between methods. However, on medium and hard tasks, our approach showed an
    evident advantage, suggesting that attention can effectively guide the search
    process in co-design. In co-design, a larger design space search volume increases
    the likelihood of discovering high-quality individuals. Notably, our method explores
    a larger search volume than other approaches. As shown in Appendix Table~\ref{tab:design_space_volume},
    AGCD achieved the largest average search volume on all tasks, followed by
    AIEA, and both are significantly larger than those of GA and BO as AGCD and AIEA
    include individuals with partial or approximate evaluation, while GA and BO
    allocate the full preset controller-training budget to every individual.
    From this perspective, AGCD
    achieves higher search efficiency. Another reason is that AGCD has higher efficiency
    in computational resource allocation. Appendix~\ref{subsubsec:fitness_distribution}
    shows the fitness distributions on selected tasks. The fitness value at the density
    peak for our method was clearly higher than those of the other baseline methods,
    and the overall distribution shifted toward higher fitness ranges.

    Figure~\ref{fig:visual} depicts visually the performance of all methods on
    the hard task GapJumper-v0. This task's objective is to jump across the gaps
    between platforms, with the reward increasng as the robot traverses more
    platforms. The results show that AGCD generated robots with crab-like running
    and jumping capabilities. In contrast, robots evolved using AIEA and BO exhibited
    unstable jumping behaviors, while those evolved using GA were obstructed at
    the gaps and failed to learn effective jumping.

    \subsection{Ablation Study}
    To validate the importance of each component of AGCD, we constructed two AGCD
    variants with all other parameters kept unchanged. The two variants are
    defined as follows:
    \begin{itemize}
        \item \textbf{AGCD-noA}: AGCD without its attention-based variation operator
            and replaced with the variation operator used in the other baseline algorithms.

        \item \textbf{AGCD-noM}: AGCD without staged evaluation or Maturity-Fitness
            Non-Dominated Sorting. Each candidate receives the full preset
            controller-training budget before selection, which uses fitness
            alone rather than fitness and training maturity.
    \end{itemize}
    \begin{figure}[t]
        \centering
        \includegraphics[width=\linewidth]{pdf/ablations.pdf}
        \caption{Mean and standard deviation of reward curves for selected tasks
        with AGCD, AGCD-noA and AGCD-noM. Each curve is averaged over five
        independent runs.}
        \label{fig:ablation}
    \end{figure}
    Figure~\ref{fig:ablation} shows the ablation study results on selected tasks.
    AGCD outperformed both AGCD-noA and AGCD-noM. On all tasks, removing a
    single component caused a sharp drop in performance, supporting the importance
    of attention-based variation and resource allocation based on training maturity
    in AGCD. The fact that AGCD outperformed AGCD-noA indicates that the
    attention-based variation operator can better focus on critical structural
    regions, thereby generating higher-quality offspring with greater potential,
    suggesting in turn that attention can guide the evolutionary search toward
    higher quality candidate solutions. In addition, AGCD outperformed AGCD-noM,
    supporting the benefit of combining staged evaluation with Maturity-Fitness
    Non-Dominated Sorting. Ranking candidates by observed fitness and training
    maturity allows promising candidates to remain in the search before their
    preset training budgets are exhausted, while elite training completes the
    budgets of selected archived candidates.

    \section{Conclusion}
    AGCD uses information from controller training in two distinct ways: its
    attention-based operators bias structural variation, and its ranking by
    fitness and training maturity distributes a limited controller-training
    budget across candidates. Training maturity records allocated training
    effort, while elite training completes the preset budgets of selected
    candidates. The
    reported Evolution Gym comparisons show stronger search efficiency and
    task performance than the evaluated baselines, while the ablations support
    the contribution of both components. These findings support
    learning-informed morphology search, although attention scores alone do not
    establish which body parts causally determine performance.

    \appendix
    % \section*{Acknowledgments}
    % This work was sponsored by the National Natural Science Foundation of China
    % (NSFC) through grant No. 52472448, and Deep Earth Probe and Mineral
    % Resources Exploration-National Science and Technology Major Project
    % throughgrants No. 2024ZD1000800 and No. 2024ZD1000804. We thank the above mentioned
    % funds for their financial support.
    %% The file named.bst is a bibliography style file for BibTeX 0.99c
    \bibliographystyle{named}
    \bibliography{ijcai26}

    \onecolumn
    % 重新定义 subsection 的编号格式为大写字母 (A, B, C...)
    \renewcommand{\thesubsection}{\Alph{subsection}}

    \section*{Appendices}
    \begin{algorithm}
        [H]
        \caption{Pseudocode of AGCD}
        \label{alg:AGCD_pseudocode}

        \textbf{Input}: Population size $N$, Maximum maturity $M_{\max}$,
        Training iterations per maturity level $T_{{m}}$,Total training
        iterations $T_{{max}}$\\
        \textbf{Output}: Optimal robot morphology $s^{*}$ and controller $\phi^{*}$

        \begin{algorithmic}
            [1]

            \item[] \textit{\# Initialization} \STATE $T_{{cur}}= 0$ \STATE $\mathcal{A}
            =0$ \textit{\# archive of robots} \STATE
            $T_{{con}}= \frac{N(N-1)}{2}\,T_{{m}}$ \textit{\# calculated
            iterations for elite training} \STATE Initialize population $P = \{P_{1}
            , P_{2}, \dots, P_{N}\}$

            \item[] \textit{\# Maturity Evolution} \WHILE{$T_{{cur}}+ T_{{con}}< T_{{max}}$}

            \FOR{each individual $p \in P$} \STATE Low-MaturityEvaluation$(p)$ \STATE
            $T_{{cur}}= T_{{cur}}+ T_{{m}}$ \STATE $p_{m}=p_{m}+1$ \STATE Update
            $\mathcal{A}$ with $p$ \ENDFOR

            \item[] \textit{\# selection} \STATE $O =$Non-DominatedSorting$(P)$ \STATE
            $S = O[1:N/2]$ \textit{\# retain top half as survivors}

            \item[] \textit{\# crossover} \STATE $C_{1}=0$ \WHILE{$|C_{1}| < N/4$}
            \STATE Randomly select two parents $p_{a}, p_{b}$ from $S$ \STATE $c
            =${AttentionCrossover}$(p_{a}, p_{b})$ \STATE Add $c$ to $C_{1}$ \ENDWHILE

            \item[] \textit{\# mutation} \STATE $C_{2}=0$ \WHILE{$|C_{2}| < N/4$}
            \STATE Randomly select one parent $p$ from $S$ \STATE $c$ =
            {AttentionMutation}$(p)$ \STATE Add $c$ to $C_{2}$ \ENDWHILE

            \item[] \textit{\# population updating} \STATE $P = S \cup C_{1}\cup
            C_{2}$

            \ENDWHILE

            \item[] \textit{\# Elite Training} \STATE $E =${Select high-potential individuals from }$\mathcal{A}$

            \FOR{each elite $e \in E$} \STATE Full-MaturityTraining$(e)$ \STATE Set
            $e_{m}= M_{\max}$ \STATE Update $\mathcal{A}$ with $e$ \ENDFOR

            \item[] \textit{\# Final Result} \STATE $(s^{*}, \phi^{*}) = \underset
            {s \in \mathcal{A}}{\arg\max}\; s_{f}$
        \end{algorithmic}
    \end{algorithm}
    \clearpage
    \begin{algorithm}
        [H]
        \caption{Maturity-Fitness Non-Dominated Sorting}
        \label{alg:mfnds} \textbf{Input}: Population $P = \{P_{1}, P_{2}, \dots,
        P_{N}\}$, where each individual $p$ has maturity $p_{m}$ and fitness
        $p_{f}$\\
        \textbf{Output}: A set of ordered frontier layers $\mathcal{F}= \{F_{1},
        F_{2}, \dots\}$
        \begin{algorithmic}
            [1] \STATE Initialize empty lists of frontiers $\mathcal{F}$ \FOR{each individual $p \in P$}
            \STATE $S_{p}\leftarrow0$ \textit{\# Set of individuals dominated by
            $p$} \STATE $n_{p}\leftarrow 0$ \textit{\# Number of individuals
            that dominate $p$} \FOR{each individual $q \in P$} \IF{$p$ dominates $q$ according to Eq.(\ref{eq:mf_dominance})}
            \STATE $S_{p}\leftarrow S_{p}\cup \{q\}$ \ELSIF{$q$ dominates $p$ according to Eq.(\ref{eq:mf_dominance})}
            \STATE $n_{p}\leftarrow n_{p}+ 1$ \ENDIF \ENDFOR \IF{$n_{p}= 0$}
            \STATE Assign $p$ to the first frontier $F_{1}$ \ENDIF \ENDFOR
            \STATE $i \leftarrow 1$ \WHILE{$F_{i}\neq0$} \STATE Initialize empty
            set $F_{i+1}$ \FOR{each individual $p \in F_{i}$} \FOR{each individual $q \in S_{p}$}
            \STATE $n_{q}\leftarrow n_{q}- 1$ \IF{$n_{q}= 0$} \STATE Assign $q$
            to frontier $F_{i+1}$ \ENDIF \ENDFOR \ENDFOR \STATE
            $i \leftarrow i + 1$ \ENDWHILE \FOR{each frontier $F_{i}$} \STATE Sort
            individuals in $F_{i}$ by descending fitness \ENDFOR \STATE \textbf{return}
            $\mathcal{F}$
        \end{algorithmic}
    \end{algorithm}

    \subsection{Method Details}
    Our overall method is summarized in Algorithm~\ref{alg:AGCD_pseudocode}. Lines
    1-4 correspond to the initialization phase, which includes population
    initialization and determining the number of iterations reserved for the
    elite training phase. Lines 5-27 describe the maturity evolution phase. At
    each generation, all individuals first receive a low-maturity evaluation,
    and the evaluated individuals are saved and updated in the robot archive. Next,
    Maturity-Fitness Non-Dominated Sorting is applied to all individuals, and the
    top half are selected as survivors. Parents are then randomly sampled from
    the survivors to perform attention-based crossover and attention-based
    mutation. Each offspring inherits the controller of the most similar parent.
    Finally, the survivors and offspring are merged to form the population for the
    next generation. Lines 28-33 correspond to the elite training phase, which aims
    to identify high-potential individuals in the robot archive that have not been
    sufficiently evaluated. High-potential individuals are selected by choosing
    the individual with the highest fitness from each of the maturity layers
    except the maximum maturity. The number of iterations reserved for elite training
    in the initialization phase is derived from this selection rule. All
    selected high-potential individuals are then trained to the maximum maturity
    level, and the robot archive is updated. Finally, the individual with the highest
    fitness is retrieved from the robot archive, and its morphology and
    controller are the output.
    \clearpage
    \subsection{Hyperparameters}
    The PPO parameters are kept the same for all tested algorithms, as listed in
    Table~\ref{tab:ppo_hyperparams}. The Transformer parameters are kept the same
    for AGCD, and variants AGCD-noA, and AGCD-noM, as listed in Table~\ref{tab:transformer_hyperparams}.
    The PPO parameters are kept the same for all tested algorithms, as listed in
    Table~\ref{tab:ppo_hyperparams}. The Transformer parameters are kept the same
    for AGCD, and variants AGCD-noA, and AGCD-noM, as listed in Table~\ref{tab:transformer_hyperparams}.
    \label{subsec:hyperparameters} \captionsetup{hypcap=false}
    \begin{table}[H]
        \centering
        \caption{Hyperparameters of PPO.}
        \label{tab:ppo_hyperparams}
        \begin{tabular}{lr}
            \toprule \textbf{Hyperparameter} & \textbf{Value}       \\
            \midrule Learning rate           & $2.5 \times 10^{-4}$ \\
            Use linear learning rate decay   & True                 \\
            Discount factor $\gamma$         & 0.99                 \\
            GAE $\lambda$                    & 0.95                 \\
            Clip parameter $\epsilon$        & 0.1                  \\
            Entropy coefficient              & 0.01                 \\
            Value loss coefficient           & 0.5                  \\
            Timesteps per rollout            & 128                  \\
            Policy epochs                    & 4                    \\
            Mini-batch size                  & 4                    \\
            \bottomrule
        \end{tabular}
    \end{table}

    \begin{table}[H]
        \centering
        \caption{Hyperparameters of the Transformer.}
        \label{tab:transformer_hyperparams}
        \begin{tabular}{lr}
            \toprule \textbf{Hyperparameter} & \textbf{Value} \\
            \midrule Number of layers        & 1              \\
            Number of attention heads        & 1              \\
            Embedding dimension              & 64             \\
            Feedforward dimension            & 128            \\
            Non linearity function           & ReLU           \\
            Dropout                          & 0.0            \\
            \bottomrule
        \end{tabular}
    \end{table}
    \clearpage
    \subsection{Details of the Environment}
    All algorithms were tested on a Linux server equipped with an Intel(R) Xeon(R)
    Gold 5220 CPU @ 2.20GHz, and no GPU was used. To accelerate training for all
    individuals in the population, we use a multithreaded parallel strategy with
    20 threads, that matches the population size. According to the settings in ~\cite{bhatia2021evolution},
    the maximum number of iterations per individual and the total number of iterations
    for each environment are listed in Table~\ref{tab:env_training_iters}.

    \begin{table}[H]
        \centering
        \small
        \begin{tabular}{cccc}
            \toprule \textbf{Env}  & \textbf{Level}          & \textbf{Individual Training Iterations} & \textbf{Total Training Iterations} \\
            \midrule Walker-v0     & \multirow{9}{*}{easy}   & 500                                     & 50000                              \\
            DownStepper-v0         &                         & 500                                     & 50000                              \\
            BeamToppler-v0         &                         & 1000                                    & 100000                             \\
            Carrier-v0             &                         & 500                                     & 50000                              \\
            Flipper-v0             &                         & 600                                     & 60000                              \\
            WingspanMazimizer-v0   &                         & 600                                     & 60000                              \\
            AreaMaximizer-v0       &                         & 600                                     & 60000                              \\
            Balancer-v0            &                         & 600                                     & 60000                              \\
            Jumper-v0              &                         & 500                                     & 50000                              \\
            \midrule Climber-v0    & \multirow{9}{*}{medium} & 400                                     & 60000                              \\
            Climber-v1             &                         & 600                                     & 90000                              \\
            Thrower-v0             &                         & 300                                     & 45000                              \\
            Pusher-v1              &                         & 600                                     & 90000                              \\
            UpStepper-v0           &                         & 600                                     & 90000                              \\
            ObstacleTraverser-v0   &                         & 1000                                    & 150000                             \\
            AreaMinimizer-v0       &                         & 600                                     & 90000                              \\
            HeightMaximizer-v0     &                         & 500                                     & 75000                              \\
            Balancer-v1            &                         & 600                                     & 90000                              \\
            \midrule BeamSlider-v0 & \multirow{9}{*}{hard}   & 1000                                    & 200000                             \\
            Carrier-v1             &                         & 1000                                    & 200000                             \\
            GapJumper-v0           &                         & 1000                                    & 200000                             \\
            Catcher-v0             &                         & 400                                     & 80000                              \\
            Lifter-v0              &                         & 300                                     & 60000                              \\
            ObstacleTraverser-v1   &                         & 1000                                    & 200000                             \\
            Climber-v2             &                         & 1000                                    & 200000                             \\
            Hurdler-v0             &                         & 1000                                    & 200000                             \\
            PlatformJumper-v0      &                         & 1000                                    & 200000                             \\
            \bottomrule
        \end{tabular}
        \caption{Training iterations across different environments.}
        \label{tab:env_training_iters}
    \end{table}
    \clearpage
    \subsection{Full Results}
    \label{subsubsec:full_results} Table~\ref{tab:final_performance} presents the
    final results of the complete comparative experiments across multiple tasks.
    Figures~\ref{fig:compare_easy},~\ref{fig:compare_medium}, and~\ref{fig:compare_hard}
    respectively display the reward curves for all easy, medium, and hard tasks.
    Both the final results in the table and the reward curves in the figures demonstrate
    that across all 27 tasks, our proposed method achieved optimal performance
    on 23 of the tasks, significantly outperforming the other baseline algorithms.
    This fully validates the effectiveness and superiority of AGCD in co-design.

    \begin{table}[H]
        \centering
        \small
        \begin{tabular}{ccccc}
            \toprule \textbf{Task} & \textbf{AGCD}        & \textbf{AIEA}      & \textbf{GA} & \textbf{BO} \\
            \midrule Walker-v0     & *\textbf{10.52±0.25} & *10.35±0.43        & 9.40±0.87   & 7.58±0.39   \\
            DownStepper-v0         & 4.66±0.26            & \textbf{4.89±0.16} & 4.18±0.56   & 4.22±0.09   \\
            BeamToppler-v0         & \textbf{2.79±0.08}   & 2.68±0.07          & 2.60±0.06   & 2.49±0.03   \\
            Carrier-v0             & \textbf{10.70±0.28}  & 4.46±0.32          & 8.71±1.75   & 2.00±0.86   \\
            Flipper-v0             & \textbf{5.09±0.30}   & 4.67±0.23          & 4.09±0.68   & 2.76±0.53   \\
            WingspanMazimizer-v0   & \textbf{0.86±0.04}   & 0.44±0.03          & 0.48±0.11   & 0.44±0.03   \\
            AreaMaximizer-v0       & \textbf{2.64±0.16}   & 1.99±0.11          & 1.45±0.11   & 1.14±0.15   \\
            Balancer-v0            & \textbf{0.12±0.01}   & 0.10±0.03          & 0.07±0.03   & 0.05±0.03   \\
            Jumper-v0              & \textbf{0.74±0.04}   & 0.58±0.04          & 0.32±0.06   & 0.28±0.03   \\
            Climber-v0             & 0.29±0.01            & \textbf{0.35±0.02} & 0.18±0.01   & 0.16±0.01   \\
            Climber-v1             & \textbf{1.35±0.07}   & 0.28±0.02          & 0.18±0.00   & 0.14±0.01   \\
            Thrower-v0             & \textbf{2.21±0.16}   & 0.94±0.06          & 0.83±0.06   & 0.84±0.02   \\
            Pusher-v1              & \textbf{0.60±0.02}   & 0.19±0.01          & 0.16±0.04   & 0.10±0.01   \\
            UpStepper-v0           & \textbf{2.54±0.13}   & 1.83±0.10          & 1.75±0.03   & 1.66±0.08   \\
            ObstacleTraverser-v0   & \textbf{5.07±0.32}   & 3.48±0.21          & 1.86±0.03   & 1.99±0.18   \\
            AreaMinimizer-v0       & \textbf{1.27±0.07}   & 0.91±0.06          & 0.71±0.03   & 0.53±0.02   \\
            HeightMaximizer-v0     & \textbf{0.36±0.02}   & 0.34±0.02          & 0.26±0.04   & 0.23±0.02   \\
            Balancer-v1            & \textbf{0.53±0.03}   & 0.41±0.02          & 0.36±0.07   & 0.40±0.01   \\
            BeamSlider-v0          & \textbf{2.83±0.17}   & 2.41±0.17          & 2.02±0.23   & 1.63±0.03   \\
            Carrier-v1             & 3.75±0.26            & \textbf{3.78±0.23} & 3.18±0.06   & 2.94±0.17   \\
            GapJumper-v0           & \textbf{8.26±0.44}   & 3.53±0.19          & 2.21±0.31   & 3.08±0.37   \\
            Catcher-v0             & \textbf{0.92±0.25}   & -2.51±0.27         & -3.59±0.26  & -3.02±0.23  \\
            Lifter-v0              & \textbf{0.10±0.03}   & 0.06±0.06          & 0.01±0.02   & -0.03±0.01  \\
            ObstacleTraverser-v1   & \textbf{5.26±0.31}   & 1.98±0.11          & 1.63±0.03   & 1.46±0.15   \\
            Climber-v2             & 0.30±0.02            & \textbf{0.42±0.03} & 0.23±0.02   & 0.30±0.01   \\
            Hurdler-v0             & \textbf{3.05±0.21}   & 1.49±0.10          & 1.28±0.08   & 1.04±0.25   \\
            PlatformJumper-v0      & \textbf{2.21±0.15}   & 1.85±0.10          & 1.75±0.03   & 1.68±0.02   \\
            \bottomrule
        \end{tabular}
        \clearpage
        \caption{Mean and standard error of all results across different tasks. Best
        results are highlighted in bold. * indicates that the result reached the
        upper bound.}
        \label{tab:final_performance}
    \end{table}
    \clearpage
    \begin{figure}[H]
        \centering
        \includegraphics[width=0.8\textwidth]{pdf/compare_easy.pdf}
        \caption{Mean and standard deviation of the reward curves on 9 easy
        tasks for AGCD, AIEA, GA, and BO. Each curve shows the mean of five
        independent runs.}
        \label{fig:compare_easy}
    \end{figure}

    \begin{figure}[H]
        \centering
        \includegraphics[width=0.8\textwidth]{pdf/compare_medium.pdf}
        \caption{Mean and standard deviation of the reward curves on 9 medium
        tasks for AGCD, AIEA, GA, and BO. Each curve shows the mean of five
        independent runs.}
        \label{fig:compare_medium}
    \end{figure}
    \clearpage
    \begin{figure}[H]
        \centering
        \includegraphics[width=0.8\textwidth]{pdf/compare_hard.pdf}
        \caption{Mean and standard deviation of the reward curves on 9 hard
        tasks for AGCD, AIEA, GA, and BO. Each curve shows the mean of five
        independent runs.}
        \label{fig:compare_hard}
    \end{figure}

    \subsection{Determination of the maximum maturity}
    \label{subsec:max_maturity}
    \begin{figure}[H]
        \centering
        \includegraphics[width=0.68\textwidth]{pdf/GapJumper-v0.pdf}
        \caption{AGCD Reward Curves for Different Maximum Maturity Levels on the
        GapJumper-v0 Task}
        \label{fig:different_maturity}
    \end{figure}

    To determine the optimal setting of the maximum maturity feature, we
    conducted a hyperparameter comparison on the GapJumper-v0 task. Specifically,
    we evaluate the performance with the maximum maturity set to 10, 20, 25, 40,
    and 50, while keeping all other experimental settings unchanged. The results
    are shown in Figure~\ref{fig:different_maturity}. When the maximum maturity was
    25, the model achieved the best overall performance on the chosen task. When
    the maximum maturity was increased or decreased from 25, the performance
    degraded to different extents. This suggests that a too small maximum maturity
    can limit search efficiency, while a too large maximum maturity can increase
    the risk of getting trapped in local optima.

    \subsection{Analysis from other perspectives}
    \subsubsection{Design space search volume}
    \label{subsubsec:search_volume} Design Space Search Volume refers to the
    total number of individuals created during the evolutionary optimization
    process, including individuals produced by operators such as crossover,
    mutation, and random generation. Each created individual is counted in this
    metric, regardless of whether it receives approximate evaluation, partial
    evaluation, or full evaluation. Table 6 reports the Design Space Search
    Volume statistics for each baseline method across all tasks. In tasks across
    all difficulty levels, the average search volume of AGCD was approximately
    twice that of AIEA, and the search volumes of both AGCD and AIEA were substantially
    higher than those of GA and BO. The latter difference mainly arises because
    all individuals generated by GA and BO require full evaluation, which is computationally
    expensive and limits the number of individuals that can be cre-ated under a fixed
    computational budget. In contrast, AGCD and AIEA introduce approximate or partial
    evaluation strategies, allowing the algorithms to create and explore more
    candidate individuals with the same resource constraint. From this perspective,
    a larger Design Space Search Volume indicates broader exploration and by
    creating and exploring substantially more individuals, AGCD increases the probability
    of discovering high-quality candidates.

    \begin{table}[H]
        \centering
        \small
        \begin{tabular}{cccccc}
            \toprule \textbf{Task}               & \textbf{Level}          & \textbf{AGCD} & \textbf{AIEA} & \textbf{GA} & \textbf{BO} \\
            \midrule Walker-v0                   & \multirow{9}{*}{easy}   & 1119          & 600           & 100         & 100         \\
            DownStepper-v0                       &                         & 1126          & 640           & 100         & 100         \\
            BeamToppler-v0                       &                         & 1130          & 600           & 100         & 100         \\
            Carrier-v0                           &                         & 1110          & 640           & 100         & 100         \\
            Flipper-v0                           &                         & 1147          & 640           & 100         & 100         \\
            WingspanMazimizer-v0                 &                         & 1110          & 600           & 100         & 100         \\
            AreaMaximizer-v0                     &                         & 1120          & 600           & 100         & 100         \\
            Balancer-v0                          &                         & 1154          & 640           & 100         & 100         \\
            Jumper-v0                            &                         & 1120          & 640           & 100         & 100         \\
            \midrule \multicolumn{2}{c}{Average} & 1126.2                  & 622.2         & 100           & 100          \\
            \midrule Climber-v0                  & \multirow{9}{*}{medium} & 1770          & 960           & 150         & 150         \\
            Climber-v1                           &                         & 1813          & 920           & 150         & 150         \\
            Thrower-v0                           &                         & 1812          & 960           & 150         & 150         \\
            Pusher-v1                            &                         & 1777          & 880           & 150         & 150         \\
            UpStepper-v0                         &                         & 1849          & 920           & 150         & 150         \\
            ObstacleTraverser-v0                 &                         & 1810          & 960           & 150         & 150         \\
            AreaMinimizer-v0                     &                         & 1764          & 960           & 150         & 150         \\
            HeightMaximizer-v0                   &                         & 1808          & 920           & 150         & 150         \\
            Balancer-v1                          &                         & 1850          & 960           & 150         & 150         \\
            \midrule \multicolumn{2}{c}{Average} & 1805.9                  & 937.8         & 150           & 150          \\
            \midrule BeamSlider-v0               & \multirow{9}{*}{hard}   & 2439          & 1160          & 200         & 200         \\
            Carrier-v1                           &                         & 2542          & 1320          & 200         & 200         \\
            GapJumper-v0                         &                         & 2380          & 1320          & 200         & 200         \\
            Catcher-v0                           &                         & 2510          & 1120          & 200         & 200         \\
            Lifter-v0                            &                         & 2490          & 920           & 200         & 200         \\
            ObstacleTraverser-v1                 &                         & 2497          & 1280          & 200         & 200         \\
            Climber-v2                           &                         & 2400          & 1320          & 200         & 200         \\
            Hurdler-v0                           &                         & 2523          & 1200          & 200         & 200         \\
            PlatformJumper-v0                    &                         & 2517          & 1280          & 200         & 200         \\
            \midrule \multicolumn{2}{c}{Average} & 2477.6                  & 1213.3        & 200           & 200          \\
            \bottomrule
        \end{tabular}
        \caption{Design space search volume across the different tasks.}
        \label{tab:design_space_volume}
    \end{table}

    \subsubsection{Fitness Distribution}
    \label{subsubsec:fitness_distribution} We present the fitness distributions
    of all methods on the Walker-v0, Carrier-v0, AreaMinimizer-v0, HeightMaximizer-v0,
    GapJumper-v0, and Lifter-v0 tasks in Figures~\ref{fig:violin_Walker-v0}-\ref{fig:violin_ObstacleTraverser-v1}.
    Overall, on most tasks, our method exhibited a higher peak fitness than the
    baseline methods, with the distribution shifted toward higher-fitness ranges.
    This suggests that our method allocates computational resources primarily to
    high-quality individuals, thereby improving the efficiency of computational resource
    allocation.
    \begin{figure}[H]
        \centering
        \includegraphics[width=0.75\textwidth]{pdf/violin_Walker-v0.pdf}
        \caption{Fitness distribution across methods on Walker-v0}
        \label{fig:violin_Walker-v0}
    \end{figure}
    \begin{figure}[H]
        \centering
        \includegraphics[width=0.75\textwidth]{pdf/violin_Carrier-v0.pdf}
        \caption{Fitness distribution across methods on Carrier-v0}
        \label{fig:violin_Carrier-v0}
    \end{figure}
    \begin{figure}[H]
        \centering
        \includegraphics[width=0.75\textwidth]{pdf/violin_AreaMinimizer-v0.pdf}
        \caption{Fitness distribution across methods on AreaMinimizer-v0}
        \label{fig:violin_Thrower-v0}
    \end{figure}
    \begin{figure}[H]
        \centering
        \includegraphics[width=0.75\textwidth]{
            pdf/violin_HeightMaximizer-v0.pdf
        }
        \caption{Fitness distribution across methods on HeightMaximizer-v0}
        \label{fig:violin_UpStepper-v0}
    \end{figure}
    \begin{figure}[H]
        \centering
        \includegraphics[width=0.75\textwidth]{pdf/violin_GapJumper-v0.pdf}
        \caption{Fitness distribution across methods on GapJumper-v0}
        \label{fig:violin_GapJumper-v0}
    \end{figure}
    \begin{figure}[H]
        \centering
        \includegraphics[width=0.75\textwidth]{pdf/violin_Lifter-v0.pdf}
        \caption{Fitness distribution across methods on Lifter-v0}
        \label{fig:violin_ObstacleTraverser-v1}
    \end{figure}
\end{document}
